% Generated draft: requires mathematical and source review.
\documentclass[11pt]{article}
\usepackage[margin=25mm]{geometry}
\usepackage{fontspec}
\setmainfont{Latin Modern Roman}
\usepackage{amsmath,amssymb,mathtools,unicode-math}
\setmathfont{Latin Modern Math}
\usepackage{enumitem,needspace,xurl,hyperref}
\setlength{\parindent}{0pt}
\setlength{\parskip}{5pt}
\setlength{\emergencystretch}{3em}
\newcommand{\N}{\mathbb{N}}
\newcommand{\Z}{\mathbb{Z}}
\newcommand{\Q}{\mathbb{Q}}
\newcommand{\R}{\mathbb{R}}
\newcommand{\C}{\mathbb{C}}
\newcommand{\F}{\mathbb{F}}
\newcommand{\A}{\mathbb{A}}
\newcommand{\PP}{\mathbb{P}}
\newcommand{\E}{\mathbb{E}}
\newcommand{\eps}{\varepsilon}
\newcommand{\dd}{\,\mathrm{d}}
\begin{document}
\section*{UnsolvedMath ID: 20001923}
\textbf{Source problem code:} catalogue entry
\subsection{Definitions and mathematical statement}
Question 8.9.1 asks: if \(M\) is a closed manifold not covered by a sphere and every geodesic of a Riemannian metric \(g\) is closed, must \((M,g)\) be a standard compact rank-one symmetric space (CROSS)? Here the Besse property means that every unit-speed geodesic is periodic, not that all prime periods must coincide. Interpret the standard CROSS list, up to scale, as round spheres, real/complex/quaternionic projective spaces and the Cayley projective plane. The supplied counterexample concerns this global list, rather than a conclusion about the universal cover or local geometry.

For every integer \(m\geq1\), on \(\C^{2m+2}\) define the antiunitary map, in pairs of coordinates, by
\[
 J(z_0,z_1,\ldots,z_{2m},z_{2m+1})
 =(-\overline z_1,\overline z_0,\ldots,
   -\overline z_{2m+1},\overline z_{2m}).
\]
Then \(J^2=-I\), so \(\tau([z])=[Jz]\) is an isometric involution of \(\C\mathrm P^{2m+1}\) with its Fubini--Study metric. It is free: a fixed point would give \(Jz=\lambda z\), hence \(-z=|\lambda|^2z\), impossible for \(z\ne0\). Put
\[
 M_m=\C\mathrm P^{2m+1}/\langle\tau\rangle
\]
with the quotient metric. This is a closed connected Besse manifold, since every geodesic lifts to a closed Fubini--Study geodesic. Its universal cover is \(\C\mathrm P^{2m+1}\), and \(\pi_1(M_m)=\Z/2\Z\).

It is not covered by a sphere: its simply connected universal cover has nonzero \(H^2\), unlike a sphere in real dimension \(4m+2\geq6\). Nor is it in the standard global CROSS list. The simply connected projective examples have the wrong fundamental group, and a real projective space has sphere universal cover. Thus this reported family refutes the literal global-CROSS implication, while each example is locally CROSS and has CROSS universal cover. It does not refute an intended conjecture with either of those weaker conclusions or an additional simply connected hypothesis.

The source also poses variants, without solving them: classify Finsler metrics on \(\C\mathrm P^n\) whose geodesics coincide with those of the standard Fubini--Study metric, and construct Lorentzian manifolds all of whose lightlike geodesics are closed. The Finsler question does not specify oriented unparameterized paths versus affine parametrizations, reversibility or normalization, so no unique classification problem is recovered beyond that wording.

The remaining extraction spills into Section 9.1 on arbitrary-signature closed-geodesic questions. It distinguishes a periodic affinely parameterized geodesic,
\[
 \gamma:\R/T\Z\to M,\qquad\nabla_{\dot\gamma}\dot\gamma=0,
\]
from a closed pregeodesic that can only locally be reparameterized affinely. In indefinite signature, a returned null velocity may be a positive nonunit multiple of the initial velocity, obstructing a global periodic affine parameter. This phenomenon cannot occur for a nonnull velocity, whose conserved nonzero squared norm forces the positive return multiplier to be one. The passage additionally chooses embeddings \(S^1\to M\), hence simple curves, as its convention; this is more restrictive than general periodic immersed-geodesic terminology and is not silently imposed on the Besse question.

The historical spillover says that many Riemannian variational methods do not extend directly to indefinite signature and reports that closed-geodesic existence in Lorentzian dimensions at least three was not understood ``to the best of our knowledge.'' This is historical source wording, not a current literature assessment. It states that a closed orientable Lorentzian surface is a torus, and attributes to [Suhr2013a] at least two simple closed geodesics on every Lorentzian two-torus, one timelike or spacelike, with examples making the count optimal. The next question introduced at the end is absent. None of these Lorentzian statements is resolved by the Fubini--Study quotient family.
\subsection{Short English statement}
Free quaternionic involutions of odd-dimensional complex projective spaces give Besse quotients that are not standard global CROSS manifolds and are not sphere-covered. They refute the literal global conclusion while preserving a CROSS universal cover. The attached Finsler and Lorentzian questions remain separate.
\subsection{Sources}
\begin{itemize}
\item Permanent UnsolvedMath ID \(20001923\), exact public code catalogue entry; local registry \texttt{aim-geometry-notes.json}, \texttt{notes[260]}, canonical tag question.
\item AIM workshop ``Geodesics,'' \emph{Open Problems and Questions about Geodesics}, extracted item 8.9, Question 8.9.1: \url{https://aimath.org/pastworkshops/geodesicsproblems.pdf}. The supplied verification note identifies the actual question on page 21 and the accidental spill into the next section. No PDF was accessed.
\item Supplied original reference labels for Lorentzian variants/background: [Mou-Suh2013], [Suhr2013b] and [Suhr2013a]. Full bibliographic details, proposer/year metadata and an original literature entry are absent.
\item Research attempt 1, ``Free Fubini--Study quotients refute the literal global-CROSS conclusion,'' labelled counterexample, refuted at the run, proved and low novelty confidence; registry solved and SOLVED-IN-LITERATURE. No full external citation is supplied for the classical free-involution construction or the assembled counterexample family.
\end{itemize}
\end{document}
